% SEGMENT OLP-0538-S01
% Part: set-theory
% Chapter: z
% Section: story

\documentclass[../../../include/open-logic-section]{subfiles}

\begin{document}

\olfileid{sth}{z}{story}
\olsection{கதையின் கூடுதல் விவரம்}

\olref[story][approach]{sec} இல், கணங்கள் உருவாகும் செயல்முறையைப் பற்றிய
ஷோன்ஃபீல்டின் விளக்கத்தை மேற்கோள் காட்டினோம். இந்தக் கதையை இன்னும் சற்றுத்
துல்லியமாக்க, மேலும் சில கொள்கைகளை இப்போது எழுத விரும்புகிறோம். அவை இதோ:
\begin{enumerate}
\item[] \stageshier. ஒவ்வொரு கணமும் ஏதோ ஒரு படிநிலையில் உருவாகிறது.
\item[] \stagesord. படிநிலைகள் வரிசைப்படுத்தப்பட்டுள்ளன: சில படிநிலைகள்
மற்றவற்றுக்கு \emph{முன்பாக} வருகின்றன.\footnote{படிநிலைகள்
\emph{நன்கு வரிசைப்படுத்தப்பட்டுள்ளன} என்று உண்மையில் நாம் மறைமுகமாக
எடுகோள்கிறோம். இதன் பொருள் என்ன என்பது \olref[ordinals][]{chap} இல்
விளக்கப்படுகிறது. இது ஒரு கணிசமான எடுகோள். உண்மையில்,
\citet{Scott1974} உருவாக்கிய மிகவும் நுட்பமான ஓர் உத்தியைப் பயன்படுத்தி,
இந்த எடுகோளை \emph{தவிர்த்து}, பின்னர் அதனை \emph{வருவிக்கலாம்}. (ஒரு
தொடக்கப் படிநிலை இருப்பதாக நாம் ஏன் கருத வேண்டும் என்பதையும் இது
விளக்கும்.) அதை இங்கு விரிவாகக் கூற முடியாது; மேலும் அறிய
\citet{ButtonLT1} ஐப் பார்க்கவும்.}
\item[] \stagesacc. எந்தப் படிநிலை $S$ க்கும், படிநிலை~$S$ க்கு
\emph{முன்பு} உருவான எந்தக் கணங்களுக்கும்: அந்தக் கணங்களைத் துல்லியமாக
உறுப்புகளாகக் கொண்ட ஒரு கணம் படிநிலை~$S$ இல் உருவாகிறது. படிநிலை~$S$
இல் வேறெதுவும் உருவாகாது.
\end{enumerate}
இவை முறைசாராக் கொள்கைகள். இருப்பினும், செர்மேலோவின் கணக் கோட்பாட்டின்
பல அடிகோள்களை ஆதரிக்க இவற்றைப் பயன்படுத்த முடியும்.

(ஓர் எச்சரிக்கையையும் வழங்க வேண்டும். முற்றிலும் தரநிலையான பெயர்களைக்
கொண்ட முற்றிலும் தரநிலையான சில அடிகோள்களை நாம் வழங்கப் போகிறோம்.
ஆனால் இப்போது சாய்வெழுத்தில் வழங்கிய கொள்கைகளுக்கு இலக்கியத்தில்
குறிப்பிட்ட பெயர்கள் எதுவும் இல்லை. உதவியாக இருக்கும் என நம்பும் இந்த
நினைவுப்பெயர்களை மட்டுமே நாம் பயன்படுத்துகிறோம்.)
% SOURCE CORRECTION TA-STH-005: supply the verb missing from “We simply monikers”.

\end{document}
